\chapter{用户使用说明书}

\providecommand{\manualscreenshot}[4][0.88\textwidth]{
\begin{figure}[H]
  \centering
  \fbox{\begin{minipage}[c][#2][c]{#1}
    \centering
    \vspace{0.25cm}
    \textbf{#3界面截图预留}\\[0.45cm]
    \small #4
  \end{minipage}}
  \caption{#3界面截图占位}
\end{figure}
}

\section{软件概述}

\subsection{软件定位}

语界 LinguaSpace 多语文旅智能导览与人才培养系统面向景区游客、外语导游专业学生、在岗导游和后台管理员提供一体化服务。系统以审核后的文旅知识库为核心，结合多语问答、语音识别、图片识别、路线推荐、AI 导游实训、真人导游协同和后台审核管理，支持游客获取可靠讲解，也支持学生和导游完成训练、接管和知识修正。

\subsection{用户角色}

\begin{longtable}{p{2.2cm}p{4.0cm}p{6.8cm}}
\toprule
角色 & 使用入口 & 主要使用内容 \\
\midrule
游客 & 游客端 H5、小程序或 Web 页面 & 文本提问、语音提问、图片问答、多语讲解、路线推荐、回答来源查看、反馈提交。 \\
学生 & 学生端 Web 页面或移动端入口 & 选择实训场景、接收虚拟游客问题、提交语音或文本讲解、查看评分报告和训练历史。 \\
导游 & 导游端 Web 页面或移动工作台 & 查看游客会话、识别高风险问题、人工接管会话、提交回答修正和协同记录。 \\
管理员 & 管理后台 Web 页面 & 用户与角色管理、知识文档上传、知识切片审核、术语维护、图谱维护、日志查询和系统状态查看。 \\
\bottomrule
\end{longtable}

\subsection{功能总览}

\begin{longtable}{p{3.0cm}p{10.4cm}}
\toprule
功能模块 & 功能说明 \\
\midrule
账户与权限 & 支持登录、退出、角色识别、权限拦截、会话过期提示和账号状态校验。 \\
智能导览 & 支持游客通过文字、语音和图片提出问题，系统基于审核知识库生成回答并展示来源。 \\
多语服务 & 支持用户选择目标语言，系统按照术语表和知识库内容生成多语讲解结果。 \\
路线推荐 & 根据游客兴趣、游览时间、当前位置和景点关系给出游览路线建议。 \\
AI 导游实训 & 为学生生成训练场景和虚拟游客问题，支持语音提交、文本提交、自动评分和报告查看。 \\
导游协同 & 为导游展示游客会话摘要、风险等级和历史问答，支持人工接管与回答修正。 \\
知识管理 & 支持管理员上传文旅资料、查看解析结果、审核知识切片、维护术语和文化图谱。 \\
日志与反馈 & 支持查询模型调用、知识检索、用户操作、异常信息和游客反馈。 \\
\bottomrule
\end{longtable}

\section{运行环境}

\subsection{服务器端环境}

\begin{longtable}{p{3.0cm}p{10.4cm}}
\toprule
类别 & 环境要求 \\
\midrule
操作系统 & 推荐使用 Linux 服务器环境；开发、演示和测试环境可使用 Windows 或 Linux 容器环境。 \\
运行时 & 后端服务运行 Python 3.10 及以上版本，前端资源由 Web 服务或静态资源服务器提供。 \\
业务服务 & 启动 API 服务、异步任务 Worker、模型适配服务、文件解析服务和定时任务服务。 \\
数据库 & PostgreSQL 保存业务数据，Milvus 或向量数据库保存知识向量，Neo4j 保存文化关系图谱。 \\
缓存与任务 & Redis 用于会话缓存、任务状态、热点数据和限流状态。 \\
对象存储 & MinIO 或兼容对象存储用于保存知识文档、图片、语音、训练音频和生成报告附件。 \\
模型服务 & ASR、翻译、视觉识别、大语言模型、向量模型和 TTS 服务可本地部署或通过模型适配器调用。 \\
网络要求 & 服务器需要能够访问数据库、缓存、对象存储和模型服务；公网部署时需配置 HTTPS。 \\
监控日志 & 服务器应开启访问日志、错误日志、模型调用日志、知识检索日志和健康检查接口。 \\
\bottomrule
\end{longtable}

\subsection{客户端环境}

\begin{longtable}{p{3.0cm}p{10.4cm}}
\toprule
客户端 & 环境要求 \\
\midrule
游客端 & 支持现代浏览器、H5 或小程序环境；语音功能需要麦克风权限，图片问答需要相机或相册权限。 \\
学生端 & 支持 Chrome、Edge、Safari 等现代浏览器；语音训练需要麦克风权限，建议使用稳定网络。 \\
导游端 & 支持移动端浏览器或 PC 浏览器；会话接管需要保持在线，建议开启新消息提醒。 \\
管理端 & 推荐使用 PC 端 Chrome 或 Edge 浏览器；知识审核、图谱维护和日志查询建议使用较大屏幕。 \\
网络条件 & 客户端需要能够访问系统服务地址；上传语音、图片和文档时需要稳定网络连接。 \\
显示要求 & 移动端适合游客和导游快速操作，PC 端适合管理员处理表格、审核和日志查询。 \\
\bottomrule
\end{longtable}

\subsection{首次使用准备}

首次使用前，系统管理员应完成服务启动、数据库初始化、基础账号创建、模型服务配置、知识库初始化和访问域名配置。普通用户只需获得系统访问地址、账号密码和对应角色权限。若使用语音、拍照或文件上传功能，客户端首次访问时需要允许浏览器或小程序申请的权限。

\section{各客户端进入方式}

\subsection{统一登录入口}

用户打开系统访问地址后进入登录页面。输入账号、密码并提交，系统根据账号角色进入游客端、学生端、导游端或管理端首页。若同一账号具备多个角色，首页显示可切换的角色入口；切换角色后，系统只展示当前角色拥有权限的菜单。

\manualscreenshot[0.86\textwidth]{5.8cm}{统一登录页}{预留登录页截图位置：系统名称、账号输入框、密码输入框、登录按钮、语言切换、忘记密码入口和错误提示区域。}

\subsection{游客端进入方式}

游客可通过景区二维码、小程序入口、H5 链接或游客端网址进入系统。免登录模式下可直接使用基础导览功能；需要保存历史记录、提交反馈或使用个性化路线时，系统提示游客登录或授权。游客端首页默认展示文本提问框、语音按钮、图片入口、语言选择和路线推荐入口。

\manualscreenshot[0.86\textwidth]{6.0cm}{游客端首页}{预留游客端首页截图位置：景区选择、语言切换、文本输入框、语音按钮、图片按钮、路线推荐入口和历史记录入口。}

\subsection{学生端进入方式}

学生使用学校或培训机构分配的账号登录。登录成功后进入实训中心首页，可查看可用训练场景、最近训练记录、待完成任务和评分报告入口。学生端入口一般位于系统首页的“AI 导游实训”或角色工作台中。

\manualscreenshot[0.86\textwidth]{6.0cm}{学生端实训首页}{预留学生实训首页截图位置：训练场景卡片、任务状态、历史记录、评分报告入口和开始训练按钮。}

\subsection{导游端进入方式}

导游使用导游账号登录后进入协同工作台。系统根据景区、班次或权限范围展示游客会话列表。高风险会话、待人工接管会话和已接管会话会按状态分组展示。导游端支持移动设备快速查看，也可在 PC 端处理较复杂的修正内容。

\manualscreenshot[0.86\textwidth]{6.0cm}{导游端协同首页}{预留导游协同首页截图位置：会话列表、风险等级、游客语言、等待时间、接管按钮和消息提醒。}

\subsection{管理端进入方式}

管理员使用后台账号登录后进入管理端首页。管理端左侧或顶部显示用户管理、知识库管理、内容审核、术语维护、图谱维护、日志查询和系统状态等菜单。无权限菜单不会显示；直接访问无权限地址时，系统返回权限不足提示。

\manualscreenshot[0.9\textwidth]{6.2cm}{管理端首页}{预留管理端首页截图位置：侧边菜单、统计卡片、待审核数量、服务状态、最近日志和快捷操作入口。}

\section{游客端详细使用说明}

\subsection{文本问答}

\begin{enumerate}[leftmargin=2em]
  \item 进入游客端首页，在景区选择区域确认当前景区。
  \item 在语言选择处选择回答语言。
  \item 在文本输入框中输入问题，例如景点历史、民族文化、参观路线、注意事项或服务设施位置。
  \item 点击发送按钮，等待系统生成回答。
  \item 查看回答内容、知识来源、可靠性提示和继续追问入口。
  \item 若回答不符合预期，可点击反馈入口，选择问题类型并提交说明。
\end{enumerate}

文本问答适合处理明确问题。若系统提示暂无可靠资料，表示当前知识库未找到审核通过的相关内容，用户可换一种问法，或提交反馈等待后台补充。

\manualscreenshot[0.86\textwidth]{6.2cm}{游客文本问答页}{预留文本问答截图位置：问题输入、语言选择、发送按钮、回答卡片、来源列表、继续追问和反馈按钮。}

\subsection{语音问答}

\begin{enumerate}[leftmargin=2em]
  \item 在游客端首页点击语音按钮。
  \item 首次使用时允许浏览器或小程序访问麦克风。
  \item 按住或点击录音控件开始说话，完成后停止录音。
  \item 系统显示语音转写文本，用户可检查转写是否正确。
  \item 确认后提交，系统按照转写文本进入问答流程。
  \item 查看回答，也可点击语音播报按钮收听讲解。
\end{enumerate}

录音时应尽量靠近麦克风并减少环境噪声。若转写结果明显错误，可重新录制或改用文字输入。

\manualscreenshot[0.86\textwidth]{6.2cm}{游客语音问答页}{预留语音问答截图位置：录音按钮、录音时长、转写文本、重新录制、提交问答、回答播报入口。}

\subsection{图片问答}

\begin{enumerate}[leftmargin=2em]
  \item 点击游客端首页的图片按钮。
  \item 选择拍照或从相册上传图片。
  \item 上传景点、文物、标识牌、菜单或服务设施图片。
  \item 根据需要补充文字问题，例如“这个建筑有什么历史”或“这道菜适合哪些忌口”。
  \item 点击提交，系统先识别图片内容，再检索知识库生成回答。
  \item 查看图片识别摘要、回答内容和知识来源。
\end{enumerate}

图片识别结果只作为检索线索，文化解释以审核知识库为准。若图片过暗、模糊或主体不清晰，系统可能提示重新上传。

\manualscreenshot[0.86\textwidth]{6.2cm}{游客图片问答页}{预留图片问答截图位置：图片预览、识别摘要、补充问题、提交按钮、回答内容和来源展示。}

\subsection{多语讲解}

\begin{enumerate}[leftmargin=2em]
  \item 在问答页或讲解页选择目标语言。
  \item 查看系统生成的目标语言讲解。
  \item 点击术语或来源入口，查看关键名词和知识出处。
  \item 如需收听讲解，点击语音播报按钮。
  \item 若译名不一致或表达不自然，可提交反馈。
\end{enumerate}

多语讲解会优先使用系统维护的术语表。景点名称、民族名称、节庆名称、非遗项目和特色饮食等专名应保持一致。

\manualscreenshot[0.86\textwidth]{6.2cm}{游客多语讲解页}{预留多语讲解截图位置：语言切换、讲解文本、术语提示、来源区域、语音播报和反馈入口。}

\subsection{路线推荐}

\begin{enumerate}[leftmargin=2em]
  \item 进入路线推荐页面。
  \item 选择游览时长、兴趣标签、同行人群和当前位置。
  \item 点击生成路线。
  \item 查看路线节点、推荐理由、预计耗时和注意事项。
  \item 根据实际情况调整兴趣或时长后重新生成。
\end{enumerate}

路线推荐为辅助建议。涉及开放时间、票务政策、临时管控和安全提示时，应以景区现场公告或系统审核知识为准。

\manualscreenshot[0.86\textwidth]{6.2cm}{游客路线推荐页}{预留路线推荐截图位置：时间选择、兴趣标签、地图区域、路线节点、推荐理由和注意事项。}

\section{学生端详细使用说明}

\subsection{选择训练场景}

\begin{enumerate}[leftmargin=2em]
  \item 登录学生端并进入实训中心。
  \item 按训练类型、目标语言、难度或景区筛选场景。
  \item 查看场景说明、训练目标、虚拟游客设定和评分维度。
  \item 点击开始训练，系统创建本次训练任务。
\end{enumerate}

训练场景包括景点讲解、民族文化介绍、接团服务、投诉处理、跨文化交流、突发情况解释和路线推荐说明等类型。

\manualscreenshot[0.86\textwidth]{6.2cm}{学生训练场景页}{预留训练场景截图位置：场景筛选、场景卡片、难度标签、虚拟游客设定、训练目标和开始按钮。}

\subsection{完成虚拟游客问答}

\begin{enumerate}[leftmargin=2em]
  \item 进入训练任务后阅读虚拟游客问题。
  \item 根据题目要求选择语音作答或文本作答。
  \item 语音作答时点击录音按钮，完成讲解后提交。
  \item 文本作答时在编辑区填写讲解稿后提交。
  \item 提交前检查是否包含事实说明、文化解释、服务礼仪和多语表达。
\end{enumerate}

学生提交后，系统会记录原始音频、转写文本、文本稿、任务上下文和评分输入，作为训练报告的依据。

\manualscreenshot[0.86\textwidth]{6.2cm}{学生讲解提交页}{预留讲解提交截图位置：虚拟游客问题、录音控件、文本稿输入区、重录按钮、提交评分按钮。}

\subsection{查看评分报告}

\begin{enumerate}[leftmargin=2em]
  \item 提交训练后等待系统生成评分报告。
  \item 查看总分、事实准确性、语言表达、跨文化礼仪、服务流程和应变能力等分项结果。
  \item 阅读系统指出的问题点和改进建议。
  \item 点击来源或证据区域，查看报告引用的知识片段或评分依据。
  \item 将报告保存到训练历史，便于后续复盘。
\end{enumerate}

评分报告用于学习反馈，不替代教师评价。若学生认为评分不准确，可提交申诉或反馈，由教师或管理员复核。

\manualscreenshot[0.86\textwidth]{6.2cm}{学生评分报告页}{预留评分报告截图位置：总分、分项评分、问题点、改进建议、评分依据和历史保存入口。}

\subsection{查看训练历史}

\begin{enumerate}[leftmargin=2em]
  \item 进入训练历史页面。
  \item 按时间、训练类型、语言、难度或得分筛选记录。
  \item 点击某条记录查看提交内容、转写文本和评分报告。
  \item 对比多次训练结果，查看薄弱维度和改进趋势。
\end{enumerate}

\manualscreenshot[0.86\textwidth]{6.2cm}{学生训练历史页}{预留训练历史截图位置：记录列表、筛选条件、得分趋势、报告查看入口和重新训练按钮。}

\section{导游端详细使用说明}

\subsection{查看游客会话}

\begin{enumerate}[leftmargin=2em]
  \item 登录导游端进入协同工作台。
  \item 在会话列表中查看游客语言、所在景区、最新问题、等待时长和风险等级。
  \item 点击会话进入详情页。
  \item 查看游客历史问题、AI 回答摘要、知识来源和系统建议。
\end{enumerate}

导游可优先处理高风险、长时间等待或多次反馈不满意的会话。会话详情帮助导游快速理解上下文，减少重复询问。

\manualscreenshot[0.86\textwidth]{6.2cm}{导游会话列表页}{预留会话列表截图位置：游客语言、景区、风险等级、最新问题、等待时长、接管按钮。}

\subsection{人工接管会话}

\begin{enumerate}[leftmargin=2em]
  \item 在会话详情页点击接管。
  \item 系统将会话状态更新为人工协同中。
  \item 导游在回复框中输入人工回复。
  \item 点击发送后，游客端收到人工回复。
  \item 处理完成后点击结束接管或转为后续跟进。
\end{enumerate}

接管适用于投诉、文化敏感问题、紧急求助、模型无法可靠回答和游客明确要求人工服务的场景。

\manualscreenshot[0.86\textwidth]{6.2cm}{导游会话接管页}{预留会话接管截图位置：会话摘要、历史问答、建议回复、人工输入框、发送按钮和结束接管。}

\subsection{提交回答修正}

\begin{enumerate}[leftmargin=2em]
  \item 在会话详情中选择需要修正的 AI 回答。
  \item 点击提交修正。
  \item 填写正确回答、修正原因、依据说明和关联景点。
  \item 提交后系统生成审核任务。
  \item 审核通过前，修正内容只作为待审核资料保存。
\end{enumerate}

导游修正用于持续改进知识库。修正内容应尽量给出明确事实依据，避免仅填写主观评价。

\manualscreenshot[0.86\textwidth]{6.2cm}{导游回答修正页}{预留回答修正截图位置：原回答、修正内容、依据说明、关联景点、提交审核按钮。}

\section{管理端详细使用说明}

\subsection{用户与角色管理}

\begin{enumerate}[leftmargin=2em]
  \item 管理员进入用户管理页面。
  \item 查看用户列表、角色、账号状态和最近登录时间。
  \item 新增用户时填写账号、姓名、角色和初始密码。
  \item 调整权限时选择角色或权限项并保存。
  \item 禁用账号时确认影响范围，保存后该账号不能继续登录。
\end{enumerate}

管理员应按照最小权限原则分配角色。学生、导游和普通管理员不应拥有超出职责范围的后台权限。

\manualscreenshot[0.9\textwidth]{6.4cm}{管理端用户权限页}{预留用户权限截图位置：用户列表、角色选择、权限矩阵、账号状态、保存按钮和登录日志。}

\subsection{知识文档上传}

\begin{enumerate}[leftmargin=2em]
  \item 进入知识库管理页面。
  \item 点击上传文档，选择文旅资料文件。
  \item 填写来源、景区、语言、版本号、主题标签和适用范围。
  \item 点击提交上传。
  \item 系统完成文件保存、文本解析、切片生成和待审核任务创建。
  \item 管理员在文档列表中查看解析状态和审核进度。
\end{enumerate}

知识资料应来自可靠来源。上传前应确认资料版权、版本、适用景区和时效性，避免将未确认内容直接用于正式问答。

\manualscreenshot[0.9\textwidth]{6.4cm}{管理端知识文档页}{预留知识文档截图位置：上传区域、文档列表、来源字段、版本号、解析状态和审核状态。}

\subsection{知识切片审核}

\begin{enumerate}[leftmargin=2em]
  \item 进入内容审核页面。
  \item 按状态、景区、来源、语言或敏感等级筛选审核任务。
  \item 打开任务详情，查看知识切片、原文来源、相似内容和系统标注。
  \item 对内容进行通过、退回、修改后通过或下线处理。
  \item 填写审核意见并提交。
  \item 审核通过的知识进入正式检索范围。
\end{enumerate}

审核时重点检查事实准确性、语言表达、文化敏感性、时效性、来源可信度和重复内容。未审核或退回内容不会进入正式回答。

\manualscreenshot[0.9\textwidth]{6.4cm}{管理端知识审核页}{预留知识审核截图位置：切片内容、来源信息、相似内容、审核意见、通过、退回和下线按钮。}

\subsection{术语维护}

\begin{enumerate}[leftmargin=2em]
  \item 进入术语维护页面。
  \item 搜索景点、民族、节庆、非遗、饮食或服务设施名称。
  \item 新增或编辑标准译名、目标语言、适用范围和说明。
  \item 保存后，系统在多语讲解和翻译结果中优先使用标准译名。
  \item 若发现冲突译名，按审核意见保留统一版本。
\end{enumerate}

术语维护直接影响多语回答一致性。管理员应避免同一专名出现多个互相矛盾的译法。

\manualscreenshot[0.9\textwidth]{6.4cm}{管理端术语维护页}{预留术语维护截图位置：术语列表、语言字段、标准译名、适用范围、冲突提示和保存按钮。}

\subsection{文化图谱维护}

\begin{enumerate}[leftmargin=2em]
  \item 进入文化图谱维护页面。
  \item 搜索或新增景点、人物、民族、节庆、非遗项目、建筑、饮食和禁忌等实体。
  \item 为实体建立关系，例如所属、相关、位于、体现、注意事项等。
  \item 查看图谱预览，确认关系方向和说明无误。
  \item 保存后，路线推荐和文化解释可读取相关关系。
\end{enumerate}

图谱关系应清晰、可解释、可追踪。错误关系可能导致路线推荐和文化讲解出现偏差。

\manualscreenshot[0.9\textwidth]{6.4cm}{管理端文化图谱页}{预留文化图谱截图位置：实体列表、关系编辑、图谱预览、关系说明和保存按钮。}

\subsection{日志查询}

\begin{enumerate}[leftmargin=2em]
  \item 进入日志查询页面。
  \item 按 requestId、用户、时间、模块、模型名称或状态筛选日志。
  \item 查看请求内容摘要、检索来源、模型调用耗时、返回状态和异常信息。
  \item 对异常问答可定位到知识片段、模型调用和用户反馈。
  \item 根据日志结果补充知识、调整术语或提交缺陷。
\end{enumerate}

日志中涉及手机号、证件号、token 等敏感字段时应脱敏展示。管理员不得将日志信息外泄。

\manualscreenshot[0.9\textwidth]{6.4cm}{管理端日志查询页}{预留日志查询截图位置：筛选条件、日志列表、requestId、模型耗时、检索来源和异常详情。}

\subsection{系统运行状态查看}

\begin{enumerate}[leftmargin=2em]
  \item 进入系统状态页面。
  \item 查看 API 服务、数据库、缓存、对象存储、模型服务和任务队列状态。
  \item 若存在异常状态，查看告警详情和最近错误日志。
  \item 处理完成后重新执行健康检查。
\end{enumerate}

系统状态页面用于日常巡检。若模型服务不可用，游客问答、图片识别、语音识别或训练评分可能受到影响。

\manualscreenshot[0.9\textwidth]{6.4cm}{管理端系统状态页}{预留系统状态截图位置：服务健康、数据库状态、模型服务状态、队列长度、告警列表和健康检查按钮。}

\section{常见操作问题}

\subsection{登录失败}

若提示用户名或密码错误，用户应确认账号、密码和大小写是否正确。若提示账号禁用或无权限访问，应联系管理员检查账号状态和角色配置。若长时间未操作后自动退出，重新登录即可继续使用。

\subsection{无法录音或拍照}

游客端或学生端无法录音时，应检查浏览器或小程序是否授权麦克风权限。无法拍照或上传图片时，应检查相机、相册和文件权限。若权限已拒绝，可在浏览器设置或系统设置中重新开启。

\subsection{问答结果为空}

系统提示暂无可靠资料时，表示当前问题未命中审核通过的知识内容。用户可更换问法、补充景区名称或提交反馈。管理员可根据反馈补充知识文档或完善术语和图谱关系。

\subsection{上传失败}

文件上传失败通常与网络中断、文件过大、格式不支持或对象存储异常有关。用户应检查文件格式和大小；管理员可查看上传日志和对象存储状态。

\subsection{训练报告未生成}

训练报告未生成时，学生可稍后刷新训练历史。若长时间仍无结果，可能是异步评分任务排队、模型服务异常或音频转写失败，管理员可在任务和模型日志中查看原因。

\section{使用注意事项}

\begin{itemize}[leftmargin=2em]
  \item 游客端回答为智能辅助讲解，涉及现场开放时间、票价、交通和安全事项时，应以景区现场公告为准。
  \item 用户上传的图片、语音和反馈应与导览或实训相关，不应上传违法、侵权或无关内容。
  \item 学生训练报告用于学习改进，不能直接替代教师评价或正式考试成绩。
  \item 导游提交修正时应填写事实依据，审核通过前不会直接进入正式知识库。
  \item 管理员审核知识时应重点关注事实准确性、来源可信度、文化敏感性和时效性。
  \item 管理员应定期检查日志、备份数据并维护模型服务、数据库、缓存和对象存储的可用性。
\end{itemize}
